\section{Radiometric framework}
\label{sec:radiometry}

\subsection{Magnitude systems and zero points}

A magnitude $m$ in a system with $F_\nu$ zero point $Z_\nu$ (in Jy)
corresponds to the monochromatic flux density
\begin{equation}
  \flam(\lambda) \;=\; Z_\nu \, \frac{c}{\lambda^2}\; 10^{-0.4\,m},
  \label{eq:flam}
\end{equation}
which the code evaluates with \code{astropy} units
(\code{magnitude\_f\_lambda}). The AB system \citep{oke1983} uses
$Z_\nu = \SI{3631}{Jy}$ at every wavelength; Vega systems use the per-filter
zero points delivered inside each SVO filter profile
\citep{rodrigo2020,bessell1998}.

\subsection{Synthetic photometry (SVO convention)}

The synthetic magnitude of a spectrum $\flam$ in a filter with response
$T_\lambda$ is
\begin{equation}
  m \;=\; -2.5\,\log_{10}
  \frac{\displaystyle\int \flam\,T_\lambda\, W_\lambda\, d\lambda}
       {\displaystyle\int \flam^{\,0}\,T_\lambda\, W_\lambda\, d\lambda},
  \qquad
  W_\lambda = \begin{cases}\lambda & \text{photon counter}\\
                            1 & \text{energy integrator,}\end{cases}
  \label{eq:synthmag}
\end{equation}
where $\flam^{\,0}$ is the zero-magnitude spectrum of
Eq.~(\ref{eq:flam}) and the weighting follows the SVO
\code{DetectorType} flag of the profile. Both conventions are implemented and
regression-tested; for a coloured source they differ measurably.

\subsection{Template calibration}
\label{sec:template}

Template spectra are absolutely calibrated CALSPEC/BPGS spectrophotometry
\citep{bohlin2014}; the catalogue field $m_{v0}$ records the visual magnitude
represented by each file. The code first restores the visual-zero
distribution, $\flam \rightarrow \flam\,10^{0.4 m_{v0}}$ (the original
Fortran convention), then scales it to the user's reference measurement:
\begin{equation}
  \flam^{\rm cal} \;=\; \flam\,10^{0.4 m_{v0}}\;
  10^{-0.4\,\left[m_{\rm ref} - (m_{\rm ref}^{\rm syn}-m_{V}^{\rm syn})\right]},
\end{equation}
where the synthetic colour term $(m_{\rm ref}^{\rm syn}-m_{V}^{\rm syn})$ is
computed from the template itself with Eq.~(\ref{eq:synthmag}) whenever the
reference filter differs from the visual response, and is exactly zero in the
original $V$-only case, which then reduces to the Fortran scale factor
$10^{-0.4(V_{\rm target}-m_{v0})}$.

\subsection{Template sources}

Templates are two-column ASCII or CALSPEC/STScI-style FITS binary tables
(\code{WAVELENGTH}/\code{FLUX}, TUNIT-aware), which admits the whole STScI
reference-atlas family: the CALSPEC standards, the solar-system
surface-brightness atlas, the galactic emission-line atlas, the transient
templates and the CLOUDY planetary-nebula grids. On import, the synthetic
Vega $V$ magnitude of each file is computed against the shipped Bessell.V
profile and stored as the catalogue $m_{v0}$; since the ETC rescales every
template to the entered target magnitude, files with arbitrary or
surface-brightness flux units remain fully usable — only the spectral shape
and the self-consistent $m_{v0}$ matter. Templates that do not fully cover
the $V$ response carry an approximate $m_{v0}$ and should be used with $V$
as the reference filter (the exact same-response path).

\subsection{Detected electron rate}

The photo-electron rate through the full system is
\begin{equation}
  \dot S \;=\; A\,\eta \int
  \flam^{\rm cal}\; T_\lambda\, Q_\lambda\, O_\lambda\,
  T_{\rm atm}(\lambda)^{X}\;
  \frac{\lambda}{hc}\; d\lambda ,
  \label{eq:rate}
\end{equation}
with the trapezoidal quadrature carried out on the filter grid
(photometry) or per resolution element (spectroscopy). Wavelength coverage is
a hard requirement: a template, QE curve, filter or atmospheric table that
does not cover the requested interval raises an error instead of silently
zero-filling.
